\chapter{Fazit und Ausblick}
In dieser Arbeit wurde das neuartige Problem DCT untersucht. 
Dabei erfolgte die Betrachtung sowohl als Entscheidungsproblem als auch Optimierungsproblem. 
Für das Entscheidungsproblem wurden zwei Modellierungen entwickelt, eine kantenbasierte und eine dreiecksbasierte Formulierung. 
Darüber hinaus wurden fünf verschiedene Instanztypen konstruiert, um sowohl lösbare als auch unlösbare Fälle zu erzeugen.
Bei der kantenbasierten Modellierung zeigte sich, dass OR-Tools in Kombination mit der Fixierung der Hülle sowie ausgewählter möglicher Kanten und unter Einbezug der alternativen Kantenbeschränkung die besten Ergebnisse erzielte. 
Mit diesem Ansatz konnten Instanzen mit einer Knotenanzahl von bis zu 220 gelöst werden. 
Für jede getestete Instanzart konnten mindestens 170 Knoten erfolgreich gelöst werden. 

Auch bei der dreiecksbasierenden Modellierung lieferte OR-Tools die besten Resultate. 
Es konnte festgestellt werden, dass diese Modellierung für alle Instanztypen, unabhängig von ihrer vermeintlichen Schwierigkeit, vergleichbare Ergebnisse erzielte. 
Zusätzlich wurde mit dem CDCL-Ansatz eine theoretische Möglichkeit aufgezeigt, den Lösungsprozess weiter zu optimieren. 
Dennoch bleibt die Skalierbarkeit auf große Instanzen eine zentrale Herausforderung, da die vorgestellten Verfahren insbesondere für kleinere Knotenmengen praktikabel sind. 
Darüber hinaus zeigte sich, dass die Struktur der Instanzen, insbesondere die Verteilung der Kantenlängen und Knotengrade, bei der kantenbasierten Formulierung einen erheblichen Einfluss auf die Schwierigkeit der Instanz hat.

Ein vielversprechender Ausblick besteht darin, die vorgestellten Verfahren im Hinblick auf ihre Skalierbarkeit zu verbessern. 
Eine interessante Möglichkeit wäre, den Arbeitsspeicherverbrauch der Ansätze systematisch zu vergleichen, da dieser insbesondere bei großen Instanzen zu Problemen geführt hat. 
Ebenso könnte die dreiecksbasierte Modellierung genauer untersucht werden, um zu verstehen, weshalb alle Instanztypen als ähnlich schwierig erschienen. 
Genauso interessant wäre es zu untersuchen, warum bei der Kantenbedingung die Optimierungen bei OR-Tools in Kombination geholfen haben und bei PySAT und Gurobi nur einzeln.
Darüber hinaus bietet der \emph{Propagation}-Ansatz weiteres Potenzial, da hierbei die Eigenschaft genutzt werden könnte, dass Kanten isoliert in Teilinstanzen auftreten. 
So können diese Kanten genutzt werden, um mit zwei weiteren die Instanz zu halbieren.
Schließlich wäre es interessant, bei der Instanzgenerierung normalverteilte Triangulierungen zu berücksichtigen, um einheitlichere und realistischere Szenarien abzubilden.

Ein interessante Abwandlung des Problems wäre es, wenn in der Triangulierung Kanten vorgegeben sind, die immer genutzt werden müssen.
Dies wurde in~\cref{sec:kanten_vorab_berechnen} teilweise betrachtet, aber nicht als eigenes Problem.
Dieses Problem könnte man noch weiter abwandeln, indem man keine Knotenmengen, sondern Polygone mit und ohne Löcher nutzt.
Auch könnte man das in~\cref{sec:einleitung:min-degree} genannte \emph{Min-Degree Constrained Minimum Spanning Tree} abwandeln, sodass der Sollgrad exakt und nicht als Mindestgrad betrachtet wird.


